\documentclass[10pt,letterpaper,oneside]{article}\usepackage{geometry}\geometry{top=1in,bottom=1in,left=1in,right=1in,includeheadfoot}\usepackage{xcolor}\definecolor{color@navy}{HTML}{1F3B63}\definecolor{color@grey}{HTML}{404040}\definecolor{color@hdr}{HTML}{F2F4F7}\usepackage{titlesec}\titleformat{\section}  {\normalfont\huge\bfseries\color@navy}  {}{0pt}{\MakeUppercase}\titleformat{\subsection}  {\normalfont\Large\bfseries\color@navy}  {}{0pt}{}\titleformat{\subsubsection}  {\normalfont\large\bfseries\color@navy}  {}{0pt}{}\titlespacing*{\section}{0pt}{14pt}{6pt}\titlespacing*{\subsection}{0pt}{10pt}{4pt}\titlespacing*{\subsubsection}{0pt}{8pt}{3pt}\setlength{\parindent}{0pt}\setlength{\parskip}{0.6em}\parindent=0pt\usepackage{enumitem}\setlist{nosep,leftmargin=*,itemsep=0pt,topsep=2pt,partopsep=0pt}\usepackage{array}\usepackage{booktabs}\usepackage{caption}\captionsetup{font={small,it},labelfont={bf},labelsep=period,  parskip=2pt,skip=4pt,list=no}\captionsetup[table]{labelfont={bf,small,color@navy},name=Table}\usepackage{graphicx}\usepackage{verbatim}\usepackage{ragged2e}\usepackage{longtable}\usepackage{tabularx}\usepackage{multirow}\usepackage{colortbl}\definecolor{rowshade}{HTML}{F2F4F7}\usepackage{tikz}\usetikzlibrary{positioning,calc}\usepackage[hidelinks]{hyperref}\hypersetup{colorlinks=true,linkcolor=color@navy,urlcolor=color@navy}\pagestyle{fancy}\fancyhf{}\renewcommand{\headrulewidth}{0.4pt}\fancyhead[L]{\small\color@navy CivicSense — Project Report}\fancyhead[R]{\small\color@grey October 11, 2026}\fancyfoot[C]{\small\color@grey An Integrated Intelligent Triage Pipelin}\fancyfoot[R]{\small\color@grey Page \thepage}\usepackage{setspace}\setstretch{1.15}\AtBeginDocument{  \hypersetup{pdftitle={An Integrated Intelligent Triage Pipeline for AI-Enhanced Citizen Complaint Management},    pdfauthor={CivicSense Project Team},    pdfsubject={Project Report},    pdfkeywords={smart cities, NLP, complaint management, e-governance},    pdfcreator={CivicSense Project Team}  }}
\begin{document}
\begin{titlepage}
\centering
\vspace*{3cm}
\begin{center}
\Huge\bfseries\color@navy CivicSense
\end{center}
\vspace{0.6cm}
\begin{center}
\huge\bfseries\color@grey Project Report
\end{center}
\vspace{0.6cm}
\centerline{==========================================}
\vspace{0.4cm}
\begin{center}
\Large\bfseries\color@navy An Integrated Intelligent Triage Pipeline for AI-Enhanced Citizen Complaint Management
\end{center}
\vspace{0.4cm}
\begin{center}
\large\itshape\color@grey A complete end-to-end implementation of the base paper: classification, prioritisation, duplicate detection, fraud scoring, incident clustering, explainable AI, dashboards and report export
\end{center}
\vfill
\begin{center}
\normalsize
Submitted in partial fulfilment of the requirements of the Major Project
\\
Prepared by: CivicSense Project Team
\\
Guided by: Project Guide, Department of Computer Engineering
\\
Date: October 11, 2026
\\
Document type: Project Report (Final Deliverable)
\\
Repository: CivicSense (Flask + civic\_sense package + SQLite demo database)
\\
Version: 2.0 (expanded edition)
\\
Classification: Internal / Academic Project Deliverable
\end{center}
\end{titlepage}
\clearpage
\phantomsection
\addcontentsline{toc}{section}{Declaration}
\section*{Declaration}

I hereby declare that the project report entitled “CivicSense: An Integrated Intelligent Triage Pipeline for AI-Enhanced Citizen Complaint Management” submitted in partial fulfilment of the requirements of the Major Project is my own work, carried out under the guidance of the project guide. This report documents the design, implementation, testing and outcomes of the project from start to finish.

All sources of data, algorithms, libraries and prior research that were consulted during the development of this project have been duly acknowledged in the References section. The demonstration dataset is synthetic and was generated specifically for this project; it does not contain any real citizen data, and the fraud and hotspot scenarios are illustrative rather than production measurements.

Every quantitative result reported in Sections 5 and 6 was measured by running the actual pipeline against a clean demonstration database, and the report itself is regenerated by a script from those measurements. Where a figure is an estimate or an illustration rather than a measurement, it is explicitly labelled as such in the text.

Signature: ______________________
Name: CivicSense Project Team
Date: October 11, 2026
\newpage
\phantomsection
\addcontentsline{toc}{section}{Certificate}
\section*{Certificate}

This is to certify that the project report entitled “CivicSense: An Integrated Intelligent Triage Pipeline for AI-Enhanced Citizen Complaint Management” is a bonafide record of work carried out by the project team in partial fulfilment of the requirements of the Major Project, under my supervision and guidance.

The work embodied in this report is original, the implementation has been executed and verified on the project’s own demonstration environment, and the report has been found satisfactory in terms of scope, depth and quality.

Guide Signature: ______________________
Designation: ______________________
Date: October 11, 2026
\newpage
\phantomsection
\addcontentsline{toc}{section}{Abstract}
\section*{Abstract}

Municipalities worldwide receive hundreds of thousands of citizen complaints every day about roads, electricity, water supply, waste management and public safety. Conventional complaint-handling systems are largely manual: they store records but do not categorise, prioritise, summarise or trust-check them, which causes delayed responses, duplicate records, misdirected resources and a rising vulnerability to coordinated spam. The base paper proposes CivicSense, an AI platform that combines Natural Language Processing (NLP) and Generative AI to automatically classify, prioritise and summarise citizen complaints, turning unstructured text into actionable governance intelligence.

This project is a complete, working implementation of that proposal. The system was extended from a linear analyser into an Integrated Intelligent Triage Pipeline with eight stages: input, preprocessing (cleaning, normalisation and locality-aware duplicate detection), classification, fraud scoring, incident clustering, weighted priority assignment, explainable-AI (XAI) generation, and dashboard/report delivery. Three research-supported enhancements were added on top of the base paper’s approach:

\item Fraud detection - a separate, explainable hybrid bad-faith risk engine that combines semantic-spam similarity, submission-frequency (flood) detection and short/generic pattern analysis.

\item Semantic incident clustering - grouping fragmented reports of the same underlying incident into locality-scoped incident clusters so that a single incident is not counted many times.

\item Explainable AI - human-readable justifications for every priority and fraud decision, so officers can audit and trust the automated triage.

The system is implemented in Python with a Flask dashboard, a REST API for complaint intake, SQLite persistence, and PDF/Excel weekly report export. The classification model achieves 82.19% accuracy (+/- 5.47%) under stratified 5-fold cross-validation on a curated corpus of 152 samples. On a clean demonstration database of 200 complaints processed end-to-end through the real pipeline, the system classified all five base-paper examples plus six additional challenge texts correctly (11 of 11), identified 23 duplicates (11%) pointing at 17 distinct originals, grouped 18 records into 11 incident clusters including one 8-complaint hotspot, and surfaced 2 High-risk and 53 Medium-risk scores. Every statistic on the dashboard reflects genuine model behaviour rather than fabricated numbers; the fraud and hotspot scenarios are synthetic demonstration data and are not production precision claims.

\textbf{Abstract Highlights}

\begin{table}[htbp]
\centering
\small
\begin{tabular}{@{}ll@{}}
\toprule
Dimension & Outcome \\
\midrule
Pipeline & Eight stages, each with an explicit input/output contract \\
Classification & 82.19\% +/- 5.47\% (stratified 5-fold CV, 152 samples) \\
Categorisation demo & 11 of 11 curated texts classified correctly \\
Priority engine & Six weighted signals with a per-decision XAI breakdown \\
Duplicate detection & 23 of 200 records flagged (11\%), 17 distinct originals \\
Incident clustering & 18 records across 11 clusters; largest cluster = 8 \\
Fraud engine & Hybrid three-signal score; 2 High-risk and 53 Medium-risk scores \\
Summarisation & Weekly digest with ward highlights and resource advice \\
Export & PDF (ReportLab) and Excel (openpyxl) on demand \\
Verification & 30 automated end-to-end checks, all passing \\
\bottomrule
\end{tabular}
\caption{Summary of the project’s headline outcomes.}
\end{table}

\emph{Keywords: Smart Cities, NLP, Complaint Management, E-Governance, Text Classification, TF-IDF, LinearSVC, Fraud Detection, Incident Clustering, Explainable AI, Flask.}
\newpage
\phantomsection
\addcontentsline{toc}{section}{Acknowledgements}
\section*{Acknowledgements}

We thank our project guide for continuous direction on scope control, for insisting that every reported metric be honestly measured rather than assumed, and for reviewing the corpus design that ultimately drove the classification accuracy from an unstable early figure to a reproducible cross-validated result.

We acknowledge the authors of the base paper, whose work defined the problem, the five complaint categories and the evaluation criteria against which this implementation is compared. We also acknowledge the open-source communities behind Flask, scikit-learn, NLTK, ReportLab, openpyxl and python-docx, on which the implementation and this document depend.

Finally, we thank the classmates and department staff who reviewed early demonstrations and pointed out the failure modes that later became Section 7 of this report. Their criticism directly improved the reliability of the final system.
\newpage
\phantomsection
\addcontentsline{toc}{section}{Table of Contents}
\section*{Table of Contents}

\addcontentsline{toc}{section}{Declaration}
\addcontentsline{toc}{section}{Certificate}
\addcontentsline{toc}{section}{Abstract}
\addcontentsline{toc}{subsection}{Abstract Highlights}
\addcontentsline{toc}{section}{Acknowledgements}
\addcontentsline{toc}{section}{Table of Contents}
\addcontentsline{toc}{section}{1. Introduction}
\addcontentsline{toc}{subsection}{1.1 Background}
\addcontentsline{toc}{subsection}{1.2 Problem Statement}
\addcontentsline{toc}{subsection}{1.3 Motivation}
\addcontentsline{toc}{subsection}{1.4 Objectives}
\addcontentsline{toc}{subsection}{1.5 Scope}
\addcontentsline{toc}{subsection}{1.6 Research Questions}
\addcontentsline{toc}{subsection}{1.7 Stakeholders and Beneficiaries}
\addcontentsline{toc}{subsection}{1.8 Expected Outcomes and Success Criteria}
\addcontentsline{toc}{subsection}{1.9 Organisation of the Report}
\addcontentsline{toc}{section}{2. Literature Review and Gap Analysis}
\addcontentsline{toc}{subsection}{2.1 Existing Complaint Platforms}
\addcontentsline{toc}{subsection}{2.2 NLP Approaches to Complaint Classification}
\addcontentsline{toc}{subsection}{2.3 Generative AI for Summarisation}
\addcontentsline{toc}{subsection}{2.4 Smart-City Adjacent Work}
\addcontentsline{toc}{subsection}{2.5 Abuse and Fraud in Public Reporting Systems}
\addcontentsline{toc}{subsection}{2.6 Duplicate Detection and Incident Clustering}
\addcontentsline{toc}{subsection}{2.7 Explainable AI in Public-Sector Decisions}
\addcontentsline{toc}{subsection}{2.8 Gap Analysis and Positioning}
\addcontentsline{toc}{subsection}{2.9 Summary of Related Work}
\addcontentsline{toc}{section}{3. Proposed System and Design}
\addcontentsline{toc}{subsection}{3.1 System Overview}
\addcontentsline{toc}{subsection}{3.2 The Eight-Stage Pipeline}
\addcontentsline{toc}{subsection}{3.3 Technology Stack and Rationale}
\addcontentsline{toc}{subsection}{3.4 Layered Architecture}
\addcontentsline{toc}{subsection}{3.5 Data Model}
\addcontentsline{toc}{subsection}{3.6 Module Design and Interfaces}
\addcontentsline{toc}{subsection}{3.7 Configuration and Tunable Parameters}
\addcontentsline{toc}{subsection}{3.8 Functional Requirements}
\addcontentsline{toc}{subsection}{3.9 Non-Functional Requirements}
\addcontentsline{toc}{subsection}{3.10 Design Decisions and Trade-offs}
\addcontentsline{toc}{subsection}{3.11 End-to-End Walkthrough of a Submission}
\addcontentsline{toc}{subsection}{3.12 Security, Privacy and Ethical Design}
\newpage
\phantomsection
\addcontentsline{toc}{section}{List of Tables}
\section*{List of Tables}

\begin{itemize}[leftmargin=1.5em, itemindent=0pt, label=\numberline{}, itemsep=2pt]
\item Summary of the project’s headline outcomes.
\item Stakeholders and how each is served by the system.
\item Pre-agreed success criteria and their measured outcomes.
\item Summary of related work and the gap each area leaves open.
\item Technology choices and the reasoning behind each.
\item The complaints table schema and the purpose of each column.
\item Modules, their responsibilities and their public functions.
\item Principal configuration parameters and their behaviour.
\item Functional requirements and their implementation status.
\item Non-functional requirements and how each is satisfied.
\item End-to-end walkthrough of one submission through the eight-stage pipeline.
\end{itemize}
\newpage
\phantomsection
\addcontentsline{toc}{section}{List of Abbreviations}
\section*{List of Abbreviations}

\begin{table}[htbp]
\centering
\small
\begin{tabular}{@{}ll@{}}
\toprule
Abbreviation & Expansion \\
\midrule
AI & Artificial Intelligence \\
API & Application Programming Interface \\
BERT & Bidirectional Encoder Representations from Transformers \\
CPGRAMS & Centralised Public Grievance Redress and Monitoring System \\
CV & Cross-Validation \\
DistilBERT & Distilled Bidirectional Encoder Representations from Transformers \\
GIS & Geographic Information System \\
IoT & Internet of Things \\
JSON & JavaScript Object Notation \\
KPI & Key Performance Indicator \\
LLM & Large Language Model \\
NLP & Natural Language Processing \\
NLTK & Natural Language Toolkit \\
PDF & Portable Document Format \\
REST & Representational State Transfer \\
SLA & Service Level Agreement \\
SQL & Structured Query Language \\
SVM & Support Vector Machine \\
TF-IDF & Term Frequency - Inverse Document Frequency \\
UI & User Interface \\
XAI & Explainable Artificial Intelligence \\
XLSX & Microsoft Excel Open XML Spreadsheet \\
\bottomrule
\end{tabular}
\end{table}
\newpage


\clearpage
\phantomsection

I hereby declare that the project report entitled “CivicSense: An Integrated Intelligent Triage Pipeline for AI-Enhanced Citizen Complaint Management” submitted in partial fulfilment of the requirements of the Major Project is my own work, carried out under the guidance of the project guide. This report documents the design, implementation, testing and outcomes of the project from start to finish.

All sources of data, algorithms, libraries and prior research that were consulted during the development of this project have been duly acknowledged in the References section. The demonstration dataset is synthetic and was generated specifically for this project; it does not contain any real citizen data, and the fraud and hotspot scenarios are illustrative rather than production measurements.

Every quantitative result reported in Sections 5 and 6 was measured by running the actual pipeline against a clean demonstration database, and the report itself is regenerated by a script from those measurements. Where a figure is an estimate or an illustration rather than a measurement, it is explicitly labelled as such in the text.

Signature: ______________________

Name: CivicSense Project Team

Date: October 11, 2026

\clearpage
\phantomsection

This is to certify that the project report entitled “CivicSense: An Integrated Intelligent Triage Pipeline for AI-Enhanced Citizen Complaint Management” is a bonafide record of work carried out by the project team in partial fulfilment of the requirements of the Major Project, under my supervision and guidance.

The work embodied in this report is original, the implementation has been executed and verified on the project’s own demonstration environment, and the report has been found satisfactory in terms of scope, depth and quality.

Guide Signature: ______________________

Designation: ______________________

Date: October 11, 2026

\clearpage
\phantomsection

Municipalities worldwide receive hundreds of thousands of citizen complaints every day about roads, electricity, water supply, waste management and public safety. Conventional complaint-handling systems are largely manual: they store records but do not categorise, prioritise, summarise or trust-check them, which causes delayed responses, duplicate records, misdirected resources and a rising vulnerability to coordinated spam. The base paper proposes CivicSense, an AI platform that combines Natural Language Processing (NLP) and Generative AI to automatically classify, prioritise and summarise citizen complaints, turning unstructured text into actionable governance intelligence.

This project is a complete, working implementation of that proposal. The system was extended from a linear analyser into an Integrated Intelligent Triage Pipeline with eight stages: input, preprocessing (cleaning, normalisation and locality-aware duplicate detection), classification, fraud scoring, incident clustering, weighted priority assignment, explainable-AI (XAI) generation, and dashboard/report delivery. Three research-supported enhancements were added on top of the base paper’s approach:

\item F r a u d   d e t e c t i o n   -   a   s e p a r a t e ,   e x p l a i n a b l e   h y b r i d   b a d - f a i t h   r i s k   e n g i n e   t h a t   c o m b i n e s   s e m a n t i c - s p a m   s i m i l a r i t y ,   s u b m i s s i o n - f r e q u e n c y   ( f l o o d )   d e t e c t i o n   a n d   s h o r t / g e n e r i c   p a t t e r n   a n a l y s i s .

\item S e m a n t i c   i n c i d e n t   c l u s t e r i n g   -   g r o u p i n g   f r a g m e n t e d   r e p o r t s   o f   t h e   s a m e   u n d e r l y i n g   i n c i d e n t   i n t o   l o c a l i t y - s c o p e d   i n c i d e n t   c l u s t e r s   s o   t h a t   a   s i n g l e   i n c i d e n t   i s   n o t   c o u n t e d   m a n y   t i m e s .

\item E x p l a i n a b l e   A I   -   h u m a n - r e a d a b l e   j u s t i f i c a t i o n s   f o r   e v e r y   p r i o r i t y   a n d   f r a u d   d e c i s i o n ,   s o   o f f i c e r s   c a n   a u d i t   a n d   t r u s t   t h e   a u t o m a t e d   t r i a g e .

The system is implemented in Python with a Flask dashboard, a REST API for complaint intake, SQLite persistence, and PDF/Excel weekly report export. The classification model achieves 82.19% accuracy (+/- 5.47%) under stratified 5-fold cross-validation on a curated corpus of 152 samples. On a clean demonstration database of 200 complaints processed end-to-end through the real pipeline, the system classified all five base-paper examples plus six additional challenge texts correctly (11 of 11), identified 23 duplicates (11%) pointing at 17 distinct originals, grouped 18 records into 11 incident clusters including one 8-complaint hotspot, and surfaced 2 High-risk and 53 Medium-risk scores. Every statistic on the dashboard reflects genuine model behaviour rather than fabricated numbers; the fraud and hotspot scenarios are synthetic demonstration data and are not production precision claims.

\subsection*{Abstract Highlights}
\begin{table}[htbp]
\centering
\small
\begin{tabular}{@{}ll@{}}
\toprule
Dimension & Outcome \\
\midrule
Pipeline & Eight stages, each with an explicit input/output contract \\
Classification & 82.19\% +/- 5.47\% (stratified 5-fold CV, 152 samples) \\
Categorisation demo & 11 of 11 curated texts classified correctly \\
Priority engine & Six weighted signals with a per-decision XAI breakdown \\
Duplicate detection & 23 of 200 records flagged (11\%), 17 distinct originals \\
Incident clustering & 18 records across 11 clusters; largest cluster = 8 \\
Fraud engine &  Hybrid three-signal score; 2 High-risk and 53 Medium-risk scores \\
Summarisation & Weekly digest with ward highlights and resource advice \\
Export & PDF (ReportLab) and Excel (openpyxl) on demand \\
Verification & 30 automated end-to-end checks, all passing \\
\bottomrule
\end{tabular}
\caption{Summary of the project’s headline outcomes.}
\end{table}

\emph{Smart Cities, NLP, Complaint Management, E-Governance, Text Classification, TF-IDF, LinearSVC, Fraud Detection, Incident Clustering, Explainable AI, Flask.}

\clearpage
\phantomsection

We thank our project guide for continuous direction on scope control, for insisting that every reported metric be honestly measured rather than assumed, and for reviewing the corpus design that ultimately drove the classification accuracy from an unstable early figure to a reproducible cross-validated result.

We acknowledge the authors of the base paper, whose work defined the problem, the five complaint categories and the evaluation criteria against which this implementation is compared. We also acknowledge the open-source communities behind Flask, scikit-learn, NLTK, ReportLab, openpyxl and python-docx, on which the implementation and this document depend.

Finally, we thank the classmates and department staff who reviewed early demonstrations and pointed out the failure modes that later became Section 7 of this report. Their criticism directly improved the reliability of the final system.

\clearpage
\phantomsection

 
Declaration
Certificate
Abstract
   Abstract Highlights
Acknowledgements
List of Tables
List of Abbreviations

1. Introduction
   1.1 Background
   1.2 Problem Statement
   1.3 Motivation
   1.4 Objectives
   1.5 Scope
   1.6 Research Questions
   1.7 Stakeholders and Beneficiaries
   1.8 Expected Outcomes and Success Criteria
   1.9 Organisation of the Report

2. Literature Review and Gap Analysis
   2.1 Existing Complaint Platforms
   2.2 NLP Approaches to Complaint Classification
   2.3 Generative AI for Summarisation
   2.4 Smart-City Adjacent Work
   2.5 Abuse and Fraud in Public Reporting Systems
   2.6 Duplicate Detection and Incident Clustering
   2.7 Explainable AI in Public-Sector Decisions
   2.8 Gap Analysis and Positioning
   2.9 Summary of Related Work

3. Proposed System and Design
   3.1 System Overview
   3.2 The Eight-Stage Pipeline
   3.3 Technology Stack and Rationale
   3.4 Layered Architecture
   3.5 Data Model
   3.6 Module Design and Interfaces
   3.7 Configuration and Tunable Parameters
   3.8 Functional Requirements
   3.9 Non-Functional Requirements
   3.10 Design Decisions and Trade-offs
   3.11 End-to-End Walkthrough of a Submission
   3.12 Security, Privacy and Ethical Design

4. Implementation
   4.1 Data Collection Layer
   4.2 Preprocessing Module
   4.3 Duplicate Detection
   4.4 Classification
   4.5 Fraud Detection
   4.6 Incident Clustering
   4.7 Weighted Priority Assignment
   4.8 Explainable AI Generation
   4.9 Summarisation and Weekly Digests
   4.10 Report Generation (PDF / Excel)
   4.11 Pipeline Orchestration
   4.12 Dashboard and User Interface
   4.13 Error Handling and Graceful Degradation
   4.14 Testing Strategy and the Smoke Suite
   4.15 Code Organisation and Repository Layout

5. Demonstration Data and Scenarios
   5.1 Demonstration Corpus
   5.2 Fraud Alert Scenario
   5.3 Incident Hotspot Scenario
   5.4 Dashboard Display Mapping
   5.5 Measured Corpus Statistics
   5.6 Reproducing the Demonstration

6. Results and Evaluation
   6.1 Classification Accuracy
   6.2 Priority Distribution
   6.3 Duplicate Detection Results
   6.4 Fraud Detection Results
   6.5 Incident Clustering Results
   6.6 Digest Generation Quality
   6.7 Automated Test Results
   6.8 System Response and Performance
   6.9 Comparison with the Base Paper
   6.10 Sample System Outputs
   6.11 Key Performance Indicators
   6.12 Threats to Validity

7. Development Challenges and Solutions
   7.1 Instability of the First Accuracy Metric
   7.2 Priority Inflation from an Overly-Broad Keyword
   7.3 Console Encoding During Testing
   7.4 Decision-Support Cards That Stayed Empty
   7.5 A Test Suite That Silently Never Ran
   7.6 A Fraud Signal That Never Fired
   7.7 Engineering Lessons

8. Limitations
9. Future Work
10. Ethical, Legal and Social Considerations
11. Project Management, Timeline and Work Distribution
12. Conclusion

References
Appendix A: How to Run the Project
Appendix B: REST API Reference
Appendix C: Source File Reference
Appendix D: Verification Checklist
Appendix E: Glossary
Appendix F: Sample Complaint Records
Appendix G: Configuration Reference
Appendix H: Sample Weekly Digest


\clearpage
\phantomsection

This section establishes why automated complaint triage is needed, states the precise problem the project solves, and defines the objectives, scope and success criteria against which the implementation is judged.

\subsection*{1.1 Background}
Urban governance depends on the timely redressal of citizen grievances. Every day, citizens lodge complaints about broken roads, power cuts, water scarcity, uncollected garbage and unsafe public spaces. These complaints arrive through many channels - mobile applications, web portals, email, call centres and social media - and in very different styles: some are formal and well structured, some use slang, and some are incomplete fragments such as “No water” or “Road bad”.

National and municipal programmes such as India’s CPGRAMS, Singapore’s OneService app and New York’s NYC 311 demonstrate the scale of the problem, but their back-ends still rely heavily on manual classification and status updates. The result, as the base paper notes, is delayed responses, duplicate records, poor categorisation and no real-time insight into where problems are concentrated. Authorities end up reacting to issues that have already become significant instead of acting on emerging trends.

Natural Language Processing (NLP) has matured to the point where text classification, similarity measurement and summarisation can be delivered cheaply and reliably on ordinary hardware. CivicSense applies these capabilities to the citizen-complaint domain, converting a raw, noisy stream of free text into a structured, prioritised and explainable body of governance intelligence.

The scale of the opportunity is easy to underestimate. A mid-sized municipal corporation with ten wards can receive several thousand complaints a month. If each complaint takes two minutes of manual reading before it can even be routed, the classification overhead alone consumes dozens of staff-hours every week before any actual repair work begins. Automating the routing decision does not eliminate the work, but it moves human attention from sorting to solving.

\subsection*{1.2 Problem Statement}
Existing complaint systems collect complaints but do not understand them. Municipal officers face four concrete problems:

\item N o   a u t o m a t i c   c a t e g o r i s a t i o n   -   a n   o f f i c e r   m u s t   r e a d   e a c h   c o m p l a i n t   m a n u a l l y   t o   d e c i d e   w h i c h   d e p a r t m e n t   i t   b e l o n g s   t o   ( r o a d s ,   e l e c t r i c i t y ,   w a t e r ,   w a s t e   o r   p u b l i c   s a f e t y ) .

\item N o   p r i o r i t i s a t i o n   -   a n   e m e r g e n c y   s u c h   a s   a   s e w a g e   o v e r f l o w   i s   t r e a t e d   t h e   s a m e   a s   a   f l i c k e r i n g   s t r e e t l i g h t .

\item N o   t r e n d   a n a l y s i s   -   p a t t e r n s   s u c h   a s   “ p o t h o l e s   r e c u r r i n g   i n   W a r d   3 ”   o r   “ g a r b a g e   d e l a y s   i n   W a r d   5 ”   a r e   i n v i s i b l e   i n s i d e   t h o u s a n d s   o f   r a w   l o g   e n t r i e s .

\item D u p l i c a t e   r e c o r d s   -   t h e   s a m e   g r i e v a n c e   i s   r e - f i l e d   a c r o s s   m u l t i p l e   c h a n n e l s ,   i n f l a t i n g   c o u n t s   a n d   w a s t i n g   e f f o r t .

CivicSense additionally recognises a fifth, emerging problem that the base paper does not address: coordinated bad-faith reporting. Because civic portals are public and easy to script, a small group - or a bot - can submit near-identical complaints to distort priority metrics and redirect resources. Any serious triage platform therefore needs a fraud-scoring stage as well as a categorisation stage.

A sixth, subtler problem is the trust deficit created by opaque automation. If a system silently promotes a complaint to “High Priority”, an officer who disagrees has no way to understand or contest the decision. CivicSense therefore treats explainability as a functional requirement rather than a nicety: every automated decision carries a human-readable justification.

\subsection*{1.3 Motivation}
The project was undertaken because the base paper describes a realistic, high-impact problem that is easy to demonstrate and verify, and because it integrates several modern AI ideas - text classification, explainable scoring and generative summarisation - into one coherent system. The subject matter of smart-city governance is a recognised national priority, and the presence of a published base paper gives the project measurable success criteria against which the implementation can be honestly compared.

A secondary motivation was engineering discipline. Many student AI projects report a single accuracy number obtained from a favourable split and stop there. This project deliberately set out to demonstrate the full lifecycle: a defined problem, a designed pipeline, a measured evaluation, a documented set of failures and their fixes, and an automated test suite that proves the system still works after every change.

A third motivation was practical usefulness. The base paper’s proposal, if realised well, could plausibly be handed to a municipal IT team and improve how they route complaints. Keeping that audience in mind shaped decisions about dependencies, offline operation and explainability: a tool that requires a GPU, an internet connection and a cloud account before it shows any value is unlikely to be adopted by a municipal corporation on a modest budget.

\subsection*{1.4 Objectives}
The specific objectives of the project are:

\item T o   b u i l d   a n   e n d - t o - e n d   c o m p l a i n t - a n a l y s i s   p l a t f o r m   i m p l e m e n t i n g   t h e   b a s e   p a p e r ’ s   p i p e l i n e   f r o m   d a t a   c o l l e c t i o n   t h r o u g h   t o   d e c i s i o n - s u p p o r t   v i s u a l i s a t i o n .

\item T o   a c h i e v e   h i g h ,   h o n e s t l y - m e a s u r e d   c l a s s i f i c a t i o n   a c c u r a c y   a c r o s s   t h e   f i v e   s e r v i c e   c a t e g o r i e s   ( R o a d s ,   E l e c t r i c i t y ,   W a t e r   S u p p l y ,   W a s t e   M a n a g e m e n t   a n d   P u b l i c   S a f e t y ) .

\item T o   a s s i g n   e x p l a i n a b l e   H i g h   /   M e d i u m   /   L o w   p r i o r i t i e s   u s i n g   t h e   s a m e   s i g n a l s   t h e   p a p e r   d e s c r i b e s :   u r g e n c y   w o r d i n g ,   c a t e g o r y   r i s k ,   e m o t i v e   l a n g u a g e ,   s a m e - w a r d   f r e q u e n c y   a n d   c o m p l a i n t   a g e .

\item T o   d e t e c t   a n d   s u p p r e s s   s e m a n t i c   d u p l i c a t e s   s o   t h a t   r e p e a t e d   r e p o r t s   o f   o n e   i s s u e   a r e   n o t   c o u n t e d   a s   i n d e p e n d e n t   i n c i d e n t s .

\item T o   a d d   a   f r a u d - d e t e c t i o n   e n g i n e   t h a t   f l a g s   c o o r d i n a t e d   o r   l o w - e f f o r t   s u b m i s s i o n s   a n d   e x p l a i n s   i t s   r e a s o n i n g .

\item T o   c l u s t e r   f r a g m e n t e d   r e p o r t s   o f   t h e   s a m e   i n c i d e n t   i n t o   h o t s p o t s   a n d   s u r f a c e   t h e m   o n   t h e   d a s h b o a r d .

\item T o   g e n e r a t e   w e e k l y   n a r r a t i v e   d i g e s t s   a n d   e x p o r t   t h e m   a s   P D F   a n d   E x c e l   r e p o r t s .

\item T o   v a l i d a t e   t h e   w h o l e   s y s t e m   w i t h   a n   a u t o m a t e d   e n d - t o - e n d   t e s t   s u i t e .

\item T o   d o c u m e n t   t h e   d e v e l o p m e n t   p r o c e s s   h o n e s t l y ,   i n c l u d i n g   t h e   d e f e c t s   d i s c o v e r e d   a n d   t h e   r e a s o n i n g   b e h i n d   e a c h   d e s i g n   d e c i s i o n .

\subsection*{1.5 Scope}
The project covers the complete software implementation: data model, NLP pipeline, trained classifier, duplicate detection, fraud scoring, incident clustering, priority engine, XAI generation, digest generator, report export, dashboard and automated tests. It uses SQLite as the central database, although the base paper suggests PostgreSQL or MongoDB for production scale, and a curated synthetic demonstration dataset that is processed through the real pipeline so that every statistic reflects genuine model behaviour.

The following are explicitly inside scope:

\item F i v e - c a t e g o r y   c o m p l a i n t   c l a s s i f i c a t i o n   w i t h   a   c o n f i d e n c e   s c o r e .

\item L o c a l i t y - a w a r e   d u p l i c a t e   d e t e c t i o n   a n d   i n c i d e n t   c l u s t e r i n g   i d e n t i f i e r s .

\item A   t h r e e - s i g n a l   f r a u d   r i s k   e n g i n e   w i t h   a n   e x p l a n a t i o n   s t r i n g   p e r   r e c o r d .

\item A   s i x - s i g n a l   w e i g h t e d   p r i o r i t y   e n g i n e   w i t h   a   n u m e r i c   b r e a k d o w n   a n d   X A I   s e n t e n c e s .

\item A   F l a s k   d a s h b o a r d   e x p o s i n g   K P I s ,   c h a r t s ,   h o t s p o t s ,   f r a u d   a l e r t s   a n d   a   k e y w o r d   c l o u d .

\item A   R E S T   i n t a k e   A P I   a n d   a   c i t i z e n - f a c i n g   w e b   f o r m .

\item W e e k l y   d i g e s t   g e n e r a t i o n   w i t h   P D F   a n d   E x c e l   e x p o r t .

\item A n   a u t o m a t e d   e n d - t o - e n d   s m o k e   s u i t e   t h a t   e x e r c i s e s   t h e   r e a l   a p p l i c a t i o n .

The following are explicitly out of scope for this version:

\item L i v e   i n t e g r a t i o n   w i t h   r e a l   m u n i c i p a l   p o r t a l s   a n d   t h e i r   a u t h e n t i c a t i o n   s y s t e m s .

\item M u l t i l i n g u a l   a n d   r e g i o n a l - l a n g u a g e   p r o c e s s i n g .

\item O f f i c e r   a u t h e n t i c a t i o n ,   r o l e - b a s e d   a c c e s s   c o n t r o l   a n d   a u d i t   t r a i l s .

\item G e o s p a t i a l   m a p   r e n d e r i n g   a n d   G P S - b a s e d   r o u t i n g   o f   f i e l d   c r e w s .

\item P r o d u c t i o n - s c a l e   s t o r a g e ,   c a c h i n g   a n d   h o r i z o n t a l   d e p l o y m e n t .

\subsection*{1.6 Research Questions}
The implementation is guided by the following research questions, which mirror and extend the base paper:

\item H o w   a c c u r a t e l y   c a n   a   l i g h t w e i g h t   T F - I D F   +   L i n e a r S V C   m o d e l   c l a s s i f y   s h o r t ,   n o i s y   m u n i c i p a l   c o m p l a i n t s   i n t o   t h e   f i v e   s e r v i c e   c a t e g o r i e s ?

\item C a n   a   d e t e r m i n i s t i c ,   m u l t i - s i g n a l   s c o r i n g   e n g i n e   p r o d u c e   p r i o r i t i e s   t h a t   m a t c h   h u m a n   i n t u i t i o n   a b o u t   u r g e n c y   w h i l e   r e m a i n i n g   f u l l y   e x p l a i n a b l e ?

\item D o e s   l o c a l i t y - a w a r e   d u p l i c a t e   d e t e c t i o n   m e a n i n g f u l l y   r e d u c e   i n f l a t e d   c o m p l a i n t   c o u n t s   w i t h o u t   s u p p r e s s i n g   g e n u i n e l y   d i s t i n c t   r e p o r t s ?

\item C a n   a   h y b r i d   o f   s e m a n t i c   s i m i l a r i t y ,   s u b m i s s i o n   f r e q u e n c y   a n d   p a t t e r n   a n a l y s i s   p r o v i d e   a   c r e d i b l e   f i r s t - p a s s   f r a u d   s i g n a l ?

\item C a n   s e m a n t i c   c l u s t e r i n g   o f   p a r a p h r a s e d   c o m p l a i n t s   r e c o v e r   a   s i n g l e   u n d e r l y i n g   i n c i d e n t   f r o m   m a n y   f r a g m e n t e d   r e p o r t s ?

\item D o e s   a   t e m p l a t e - b a s e d   d e t e r m i n i s t i c   s u m m a r i s e r   p r o d u c e   d i g e s t s   o f   a   q u a l i t y   c o m p a r a b l e   t o   t h e   b a s e   p a p e r ’ s   i l l u s t r a t i v e   e x a m p l e s   w i t h o u t   r e q u i r i n g   a   l a n g u a g e - m o d e l   A P I ?

\subsection*{1.7 Stakeholders and Beneficiaries}
The system has several distinct user groups, and the design was shaped by their differing needs:

\begin{table}[htbp]
\centering
\small
\begin{tabular}{@{}lll@{}}
\toprule
Stakeholder & Primary need & How CivicSense serves them \\
\midrule
Citizens & Fast, transparent acknowledgement that a complaint was received and routed & Web form with instant AI analysis; status tracking; API intake from partner apps \\
Front-line officers & A trustworthy queue ordered by real urgency & Explainable High/Medium/Low priority with a numeric breakdown per complaint \\
Department heads & Evidence of where resources are needed & Ward hotspot bars, category mix, incident clusters and resource-allocation advice \\
Municipal administrators & Periodic, readable situational awareness & Weekly narrative digest exportable as PDF and Excel \\
Data / integrity teams & Early warning of abuse or automation & Fraud Alerts card with scores, levels and contributing signals \\
Auditors and oversight & A defensible explanation for every automated decision & XAI sentences stored with each record for priority and fraud \\
\bottomrule
\end{tabular}
\caption{Stakeholders and how each is served by the system.}
\end{table}

\subsection*{1.8 Expected Outcomes and Success Criteria}
The project was judged against explicit, pre-agreed criteria rather than a general impression. These criteria were fixed before the final evaluation run:

\begin{table}[htbp]
\centering
\small
\begin{tabular}{@{}lll@{}}
\toprule
Criterion & Target & Measured outcome \\
\midrule
Classification accuracy & Above 80\% under cross-validation & 82.19\% +/- 5.47\% \\
Base-paper examples & All five classified correctly & 5 of 5 correct \\
Priority explainability & Every decision carries a reason & Provided for every record \\
Duplicate handling & Flag re-files without over-flagging & 23 of 200 flagged (11\%) \\
Fraud visibility & At least one High-risk row visible on the dashboard & 2 High and 53 Medium scores \\
Hotspot visibility & A populated cluster on the dashboard & 8-complaint cluster in Wakad \\
Report export & PDF and Excel both downloadable & Both verified by the smoke suite \\
Automated verification & A repeatable end-to-end suite & 30 checks, all passing \\
\bottomrule
\end{tabular}
\caption{Pre-agreed success criteria and their measured outcomes.}
\end{table}

\subsection*{1.9 Organisation of the Report}
The remainder of the report is organised as follows. Section 2 reviews related work and identifies the gaps that CivicSense addresses. Section 3 presents the proposed system and its design, including functional and non-functional requirements. Section 4 documents the implementation of each module in detail. Section 5 describes the demonstration data and the two explicit scenarios. Section 6 reports the results and evaluates them against the base paper. Section 7 records the development challenges and how they were solved. Section 8 discusses limitations. Section 9 lays out future work. Section 10 covers ethical, legal and social considerations. Section 11 describes project management and the timeline. Section 12 concludes. References and appendices A to H follow.

\clearpage
\phantomsection

This section summarises the landscape the base paper reviews, adds the adjacent literature on abuse detection, duplicate handling and explainability that the extension builds on, and positions the project within it.

\subsection*{2.1 Existing Complaint Platforms}
CPGRAMS (India), OneService (Singapore) and NYC 311 (United States) are the best-known large-scale citizen-grievance platforms. They improved transparency and accessibility by giving citizens a single submission point and a tracking number, and they generate large volumes of structured status data. What they did not solve is the analytical problem: categorisation, prioritisation and trend discovery remain largely manual, and the systems do not attempt to detect deliberate abuse of the intake channel.

A useful way to think about these platforms is that they are excellent record systems and poor decision systems. They answer “what was submitted?” reliably, but they do not answer “what matters most right now?” without a human reading the queue. CivicSense is deliberately a decision layer rather than another record layer: it assumes complaints arrive from somewhere (a portal, an app, an email gateway) and concentrates on what happens next.

\subsection*{2.2 NLP Approaches to Complaint Classification}
Early research on municipal complaint classification used text mining with Naive Bayes and Support Vector Machines. These approaches work reasonably well on clean text but struggle with multilingual, ambiguous or informal complaints. The introduction of transformer models such as BERT, RoBERTa and DistilBERT pushed classification accuracy past 90% on several municipal datasets, at the cost of substantially larger compute and memory requirements.

A middle path is to combine TF-IDF feature extraction with a linear classifier. This delivers strong accuracy on short, keyword-heavy texts such as complaints, trains in seconds on a laptop, needs no GPU, and is fully interpretable. CivicSense adopts this as the primary model and keeps a zero-shot DistilBERT path available behind a feature flag for demonstration purposes.

The interpretability argument deserves emphasis. A linear SVM’s decision is a weighted sum of term features, so it is straightforward to inspect which words drive a particular classification. A transformer’s decision is not directly inspectable. For a public-sector triage tool where an officer may have to defend an automated routing decision, the simpler model has a real advantage beyond its speed.

\subsection*{2.3 Generative AI for Summarisation}
Generative models have been used in healthcare and legal services to produce patient summaries and case briefs. The base paper observes that this capability has not been widely applied to citizen-to-government complaint analysis. A city administrator does not want a list of one thousand complaints; they want a weekly digest of the form: “Ward 7 had 35 road complaints this week, the recurring concern being potholes - High Priority.” Closing this gap is the central motivation of CivicSense.

Two implementation strategies exist for this kind of narrative output. The first is to call a large language model with the period’s statistics and let it write the prose. The second is to compute the same statistics deterministically and generate the narrative from a template. The first produces more fluent language but depends on network access, cost and an API key. The second is reproducible, free and always available, at the cost of some fluency. CivicSense implements both and prefers the template engine by default, so the demonstration never depends on an external service.

\subsection*{2.4 Smart-City Adjacent Work}
Predictive models forecast water-pipe bursts, machine learning optimises traffic flow, and IoT-connected garbage bins streamline collection routes. These systems optimise municipal assets from historical telemetry. Complaint analysis optimises a different signal - the citizen’s voice - and feeds it into the same decisions. Sentiment analysis of social-media posts measures public mood but cannot pinpoint a specific service failure at ward level. CivicSense combines classification, prioritisation and summarisation to close that gap.

The distinction matters practically. An IoT bin sensor tells you a bin is full; it does not tell you that forty citizens have already complained about overflowing bins on the same street, which is a different and stronger signal about service failure. Combining citizen text with asset telemetry is a natural future step, and is listed in Section 9.

\subsection*{2.5 Abuse and Fraud in Public Reporting Systems}
Open reporting channels are inherently vulnerable to abuse. A body of work on spam and Sybil attacks in public platforms has established three recurring patterns relevant to civic complaint intake: bursts of activity concentrated in a narrow time window, near-identical content submitted repeatedly, and content that is too short or generic to be verified. These are the same signals that fraud engines use in review and social platforms, and they transfer naturally to the complaint domain.

The design consequence is that fraud detection should be kept separate from deduplication. A genuine citizen re-filing a legitimate complaint and a bot flooding a channel can produce the same text and, superficially, the same similarity score. Duplicate detection asks “is this the same issue?” whereas fraud detection asks “was this filed in bad faith?” Merging them would punish legitimate duplicates. CivicSense therefore implements them as two independent modules with distinct outputs.

\subsection*{2.6 Duplicate Detection and Incident Clustering}
Classical deduplication compares normalised text with a similarity measure such as Jaccard overlap, cosine similarity over TF-IDF vectors, or edit-distance ratios. These approaches are cheap and deterministic but brittle to paraphrase: “a tree fell on the road” and “a fallen branch blocked traffic” share almost no surface tokens while describing the same incident.

Embedding-based approaches solve this by comparing meaning rather than wording. Sentence-BERT and similar encoders map paraphrases to nearby vectors, so a single similarity threshold can group them. Incident clustering is the natural extension of this idea: rather than merely flagging one record as a duplicate of another, reports are grouped into clusters that represent distinct open incidents. This shift from pairwise deduplication to clustering is what allows a system to answer “how many distinct incidents are open?” instead of “how many records exist?”.

A practical difficulty is that embedding libraries add heavy dependencies. CivicSense addresses this by making embeddings the preferred path and providing a deterministic difflib-based fallback on normalised text, so the clustering feature degrades gracefully instead of disappearing when the embedding stack is unavailable.

\subsection*{2.7 Explainable AI in Public-Sector Decisions}
Automated decisions in public administration carry a legitimacy requirement that decisions in many other domains do not. A citizen whose complaint is deprioritised is entitled to know why, and an officer who overrides a machine ranking needs to understand what the machine thought it saw. Explainable AI research distinguishes between intrinsic explainability (the model is simple enough to inspect) and post-hoc explainability (a separate method explains the model).

CivicSense uses both. The classifier is intrinsically explainable because it is a linear model over term features. The priority and fraud engines, which combine several weighted signals, are explained post-hoc by converting their numeric breakdowns into natural-language sentences. This is a lighter-weight and more auditable approach than attributing a deep model’s output, and it is well matched to the domain, where the contributing factors are already defined and bounded.

\subsection*{2.8 Gap Analysis and Positioning}
Synthesising the above, the base paper establishes a preprocessing-plus-classification pipeline but leaves four gaps that this project addresses:

\item \textbf{Duplicate handling: } D e d u p l i c a t i o n   i s   n a i v e   a n d   d o e s   n o t   a c c o u n t   f o r   l o c a l i t y   o r   p a r a p h r a s i n g .   C i v i c S e n s e   m a k e s   i t   l o c a l i t y - a w a r e   a n d   s i m i l a r i t y - b a s e d .

\item \textbf{Fraud handling: } B a d - f a i t h   c o n t e n t   i s   n o t   m o d e l l e d ,   s o   c o o r d i n a t e d   s p a m   c a n   s i l e n t l y   d i s t o r t   p r i o r i t y   a n d   r e s o u r c e   a l l o c a t i o n .   C i v i c S e n s e   a d d s   a   s e p a r a t e ,   e x p l a i n a b l e   f r a u d - s c o r i n g   e n g i n e .

\item \textbf{Incident handling: } I n d i v i d u a l   c o m p l a i n t s   a r e   t r a c k e d ,   b u t   f r a g m e n t e d   r e p o r t s   o f   o n e   i n c i d e n t   a r e   n o t   g r o u p e d .   C i v i c S e n s e   c l u s t e r s   t h e m   i n t o   i n c i d e n t s .

\item \textbf{Explainability: } T r i a g e   s c o r e s   a r e   p r e s e n t e d   w i t h o u t   j u s t i f i c a t i o n ,   w h i c h   r e d u c e s   o f f i c e r   t r u s t   a n d   a d o p t i o n .   C i v i c S e n s e   e m i t s   a n   X A I   e x p l a n a t i o n   f o r   e v e r y   p r i o r i t y   a n d   f r a u d   d e c i s i o n .

The design follows the software-engineering principle of separation of concerns: each stage is a discrete module with a well-defined input and output contract, so each can be tested, replaced or extended independently.

\subsection*{2.9 Summary of Related Work}
\begin{table}[htbp]
\centering
\small
\begin{tabular}{@{}llll@{}}
\toprule
Area & Representative approach & Strength & Limitation addressed by CivicSense \\
\midrule
Complaint platforms & CPGRAMS, OneService, NYC 311 & Scale, tracking, transparency & No automated triage or abuse detection \\
Classical text classification & Naive Bayes, SVM over bag-of-words & Cheap, interpretable & Weak on informal and ambiguous text \\
Transformer classification & BERT, RoBERTa, DistilBERT & High accuracy & Heavy dependencies; opaque decisions \\
Deduplication & Jaccard / cosine / edit distance & Deterministic, fast & Brittle to paraphrase; no locality awareness \\
Embedding similarity & Sentence-BERT and similar encoders & Handles paraphrase & Heavy dependency; needs a fallback \\
Incident grouping & Clustering over embeddings & Counts distinct incidents & Rarely wired into a live triage dashboard \\
Abuse detection & Burst, repetition and generic-text signals & Catches coordinated spam & Usually absent from civic complaint systems \\
Summarisation & Extractive scoring; LLM narrative & Produces readable digests & LLM path depends on external services \\
Explainability & Post-hoc justification from score breakdowns & Builds officer trust & Often omitted from triage tools \\
\bottomrule
\end{tabular}
\caption{Summary of related work and the gap each area leaves open.}
\end{table}

\clearpage
\phantomsection

This section presents the system as designed: its overall shape, the pipeline it implements, the technologies chosen and why, the data model, the module contracts, the requirements it must satisfy and the trade-offs made along the way.

\subsection*{3.1 System Overview}
CivicSense is delivered as a Flask web application backed by a SQLite database and a Python package, civic_sense, that contains all analytical logic. Citizens submit complaints through a web form or a REST endpoint; every submission is processed synchronously by the pipeline before it is stored, so no unanalysed record ever enters the database. Officers interact with a single-page dashboard that surfaces key performance indicators, the latest digest, category and priority mixes, ward hotspots, trending keywords and recent complaints.

The choice to analyse before storing is deliberate and has several consequences. It guarantees that a stored record always carries a complete set of derived fields, which means dashboards and reports never have to handle partially analysed rows. It also means a classification or clustering failure surfaces immediately at submission time rather than silently corrupting later aggregates. The cost is that submission latency includes analysis time, which is acceptable at demonstration scale and is measured in Section 6.8.

\subsection*{3.2 The Eight-Stage Pipeline}
The implemented flow is:

Input -> Preprocessing (cleaning + normalisation + locality-aware duplicate detection) -> Classification -> Fraud Scoring -> Incident Clustering -> Weighted Priority (with optional cluster boost) -> XAI Generation -> Dashboard / Reports

Each stage contributes a distinct artifact to the record: normalised text and a duplicate flag, a category with a confidence score, a fraud score with a level and explanation, a cluster identifier, a priority label with a numeric breakdown, and finally the stored record that drives the dashboard and reports.

\begin{figure}[htbp]
\centering
\begin{verbatim}
[1] Input            -> raw text, ward, source
        |
[2] Preprocessing    -> normalised text, duplicate flag
        |
[3] Classification   -> category, confidence
        |
[4] Fraud Scoring    -> score, level, explanation
        |
[5] Clustering       -> cluster_id
        |
[6] Priority         -> label, numeric breakdown
        |
[7] XAI Generation   -> justification sentences
        |
[8] Delivery         -> dashboard, API, reports
\end{verbatim}
\caption{The eight-stage triage pipeline and the artifact each stage produces.}
\end{figure}

\subsection*{3.3 Technology Stack and Rationale}
\begin{table}[htbp]
\centering
\small
\begin{tabular}{@{}lll@{}}
\toprule
Component & Choice & Why \\
\midrule
Language & Python 3.11/3.12 & Rich NLP and ML ecosystem; fast to develop and easy to demonstrate \\
Web framework & Flask & Minimal overhead; the project is a demonstration of the AI pipeline, not a heavy web product \\
Classifier & TF-IDF + LinearSVC (scikit-learn) & Strong on short keyword-heavy text; trains in seconds; no GPU; explainable \\
Optional classifier & Zero-shot DistilBERT & Demonstrates the paper’s BERT-family approach when the transformer stack is available \\
Similarity & Sentence-Transformers, difflib fallback & Semantic similarity when available; deterministic fallback otherwise \\
Persistence & SQLite & Zero-configuration demo database; PostgreSQL/MongoDB intended for production \\
Reporting & ReportLab + openpyxl & PDF and Excel export exactly as the paper requires \\
Document generation & python-docx & This report is itself generated reproducibly from measured data \\
UI & Hand-written HTML/CSS & Fully self-contained; the demo runs entirely offline \\
\bottomrule
\end{tabular}
\caption{Technology choices and the reasoning behind each.}
\end{table}

\subsection*{3.4 Layered Architecture}
The system is organised into four layers with strict downward dependencies:

\item P r e s e n t a t i o n   l a y e r   -   t e m p l a t e s /   a n d   s t a t i c / c s s / s t y l e . c s s :   d a s h b o a r d ,   c o m p l a i n t   l i s t ,   i n t a k e   f o r m ,   r e p o r t   v i e w s   a n d   p i p e l i n e   e x p l a i n e r .

\item A p p l i c a t i o n   l a y e r   -   a p p . p y :   F l a s k   r o u t e s   f o r   t h e   w e b   U I   a n d   R E S T   A P I ,   p l u s   t h e   r u n \_ p i p e l i n e ( )   o r c h e s t r a t o r   t h a t   c h a i n s   t h e   s t a g e s   f o r   e v e r y   i n c o m i n g   c o m p l a i n t .

\item C o r e   l a y e r   -   t h e   c i v i c \_ s e n s e   p a c k a g e :   p r e p r o c e s s i n g ,   c l a s s i f i e r ,   f r a u d \_ d e t e c t o r ,   p r i o r i t i z e r ,   s u m m a r i z e r ,   r e p o r t s ,   d a t a b a s e ,   c o n f i g ,   s e e d \_ d a t a .

\item D a t a   l a y e r   -   d a t a / c i v i c s e n s e . d b   ( c o m p l a i n t s   a n d   r e p o r t s   t a b l e s )   a n d   m o d e l s /   ( t h e   t r a i n e d   c l a s s i f i e r   a n d   i t s   m e t r i c s   J S O N ) .

The core layer never imports Flask and the presentation layer never performs analysis. This separation means the analytical pipeline can be exercised directly from a script or a test without starting a web server - which is what the smoke suite and the reporting script both rely on.

\subsection*{3.5 Data Model}
The complaints table persists every field the pipeline produces:

\item C o r e   f i e l d s :   i d ,   w a r d ,   c a t e g o r y ,   t i t l e ,   d e s c r i p t i o n ,   s o u r c e ,   c r e a t e d \_ a t ,   s t a t u s   a n d   p r i o r i t y

\item D u p l i c a t e   f i e l d s :   i s \_ d u p l i c a t e   a n d   d u p l i c a t e \_ o f

\item F r a u d   f i e l d s :   f r a u d \_ s c o r e ,   f r a u d \_ l e v e l   a n d   f r a u d \_ e x p l a n a t i o n

\item C l u s t e r i n g :   c l u s t e r \_ i d

\item P r i o r i t y   a n d   X A I :   p r i o r i t y \_ s c o r e ,   p r i o r i t y \_ b r e a k d o w n   a n d   x a i \_ e x p l a n a t i o n

A second reports table stores generated weekly digests together with their statistics, so that exported PDF and Excel files can be regenerated on demand.

\begin{table}[htbp]
\centering
\small
\begin{tabular}{@{}lll@{}}
\toprule
Column & Type & Purpose \\
\midrule
id & INTEGER PK & Surrogate key used in URLs and API responses \\
ward & TEXT & Locality; scopes duplicate detection and frequency scoring \\
category & TEXT & One of the five service categories assigned by the classifier \\
title & TEXT & Short human-readable label shown in tables \\
description & TEXT & The raw complaint text submitted by the citizen \\
source & TEXT & Intake channel (Mobile App, Web Portal, Email, Call Center, Social Media) \\
priority & TEXT & High / Medium / Low label from the priority engine \\
priority\_score & REAL & Numeric urgency score before thresholding \\
priority\_breakdown & TEXT (JSON) & Per-signal contribution for XAI \\
status & TEXT & Open / In Progress / Resolved work state \\
confidence & REAL & Classifier pseudo-confidence in [0, 1] \\
is\_duplicate & INTEGER & 1 when the record matches an earlier complaint \\
duplicate\_of & INTEGER & id of the original complaint when is\_duplicate = 1 \\
fraud\_score & REAL & 0-100 bad-faith risk score \\
fraud\_level & TEXT & High / Medium / Low risk band \\
fraud\_explanation & TEXT (JSON) & Contributing fraud signals as sentences \\
cluster\_id & INTEGER & Shared identifier for one underlying incident \\
xai\_explanation & TEXT (JSON) & Priority justification sentences \\
created\_at & TEXT & ISO timestamp used by the age and frequency signals \\
\bottomrule
\end{tabular}
\caption{The complaints table schema and the purpose of each column.}
\end{table}

\subsection*{3.6 Module Design and Interfaces}
Each analytical concern lives in its own module with a narrow interface:

\begin{table}[htbp]
\centering
\small
\begin{tabular}{@{}lll@{}}
\toprule
Stage & Module & Key functions \\
\midrule
Preprocessing & civic\_sense.preprocessing & clean\_text, tokenize, lemmatize, remove\_stopwords, normalize, similarity, find\_duplicate, detect\_incident\_cluster \\
Classification & civic\_sense.classifier & load\_classifier, classify\_text, train, get\_metrics \\
Fraud scoring & civic\_sense.fraud\_detector & analyze\_fraud\_risk, generate\_fraud\_explanation \\
Incident clustering & civic\_sense.preprocessing & detect\_incident\_cluster \\
Priority & civic\_sense.prioritizer & assign\_priority, priority\_explanation \\
Seeding & civic\_sense.seed\_data & seed\_if\_empty, \_fraud\_scenario, \_hotspot\_scenario \\
Persistence & civic\_sense.database & insert\_complaint, get\_complaint, get\_complaints, get\_stats, fraud\_alerts, cluster\_summary \\
Summarisation & civic\_sense.summarizer & extractive\_summary, generate\_weekly\_report \\
Export & civic\_sense.reports & generate\_pdf, generate\_excel \\
App routes & app.py & index, complaints, new\_complaint, api\_submit, report\_view, report\_pdf, report\_xlsx \\
\bottomrule
\end{tabular}
\caption{Modules, their responsibilities and their public functions.}
\end{table}

\subsection*{3.7 Configuration and Tunable Parameters}
All thresholds and scenario knobs live in civic_sense/config.py and the relevant modules, so that behaviour can be tuned without editing algorithm code. The principal parameters are:

\begin{table}[htbp]
\centering
\small
\begin{tabular}{@{}lll@{}}
\toprule
Parameter & Value & Effect \\
\midrule
SIM\_THRESHOLD & 0.85 & Similarity above which a complaint counts as a semantic duplicate. \\
FREQ\_THRESHOLD & 5 & Same-ward submissions in a short window that trigger the flood signal. \\
HIGH\_FRAUD\_SCORE & 70 & Fraud score at or above which the level is High. \\
MEDIUM\_FRAUD\_SCORE & 30 & Fraud score at or above which the level is Medium. \\
Priority thresholds & 3.2 / 1.6 & Score boundaries for High and Medium priority. \\
FRAUD\_SCENARIO\_COUNT & 8 & Number of seeded complaints in the fraud-alert scenario. \\
FRAUD\_SCENARIO\_WARD & Hinjewadi & Ward used by the fraud-alert scenario. \\
HOTSPOT\_SCENARIO\_WARD & Wakad & Ward used by the incident-hotspot scenario. \\
HOTSPOT\_SCENARIO\_CATEGORY & Public Safety & Category used by the incident-hotspot scenario. \\
USE\_TRANSFORMERS & 0 (default) & Enables the optional DistilBERT classification path. \\
OPENAI\_API\_KEY & unset & Enables language-model-written weekly digests. \\
\bottomrule
\end{tabular}
\caption{Principal configuration parameters and their behaviour.}
\end{table}

Because these values are explicit and centralised, the demonstration scenarios can be made more or less aggressive without touching the pipeline itself - an important property for a system whose behaviour must be explained and audited.

A deliberate constraint is that no parameter is hidden inside algorithm code. An early version of the project embedded an urgency keyword list directly inside the scoring function, which made the priority bug described in Section 7.2 difficult to reason about. Moving all such lists and thresholds to named constants made the behaviour inspectable and the fix obvious.

\subsection*{3.8 Functional Requirements}
The functional requirements were derived from the base paper and the stakeholder analysis:

\begin{table}[htbp]
\centering
\small
\begin{tabular}{@{}lll@{}}
\toprule
ID & Requirement & Status \\
\midrule
FR-1 & Accept complaints from a web form and a REST API & Implemented \\
FR-2 & Record the intake channel (source) of each complaint & Implemented \\
FR-3 & Normalise complaint text deterministically & Implemented \\
FR-4 & Classify each complaint into one of five categories & Implemented \\
FR-5 & Report a confidence score with each classification & Implemented \\
FR-6 & Detect and record semantic duplicates within a ward & Implemented \\
FR-7 & Score each complaint for bad-faith risk and record the level & Implemented \\
FR-8 & Explain every fraud score in human-readable terms & Implemented \\
FR-9 & Group paraphrased reports into incident clusters & Implemented \\
FR-10 & Assign an explainable High/Medium/Low priority & Implemented \\
FR-11 & Expose KPIs, charts and hotspots on a dashboard & Implemented \\
FR-12 & Generate weekly narrative digests & Implemented \\
FR-13 & Export digests as PDF and Excel & Implemented \\
FR-14 & Allow officers to update complaint status & Implemented \\
FR-15 & Provide an automated end-to-end test suite & Implemented \\
\bottomrule
\end{tabular}
\caption{Functional requirements and their implementation status.}
\end{table}

\subsection*{3.9 Non-Functional Requirements}
\begin{table}[htbp]
\centering
\small
\begin{tabular}{@{}lll@{}}
\toprule
ID & Requirement & How it is met \\
\midrule
NFR-1 & Offline operation with no mandatory external service & All analysis is local; the LLM digest path is optional \\
NFR-2 & Deterministic, reproducible analysis & Normalisation is stateless; seeds are fixed; metrics are stored \\
NFR-3 & Graceful degradation when optional dependencies are missing & Difflib fallback, sklearn fallback, template digest fallback \\
NFR-4 & Explainability of every automated decision & Numeric breakdowns are converted to sentences and stored \\
NFR-5 & Low setup cost for demonstration & A single run script installs dependencies and starts the app \\
NFR-6 & No real personal data in the repository & All demonstration records are synthetic \\
NFR-7 & Reasonable responsiveness at demonstration scale & Full pipeline per submission; measured in Section 6.8 \\
NFR-8 & Portability across platforms & Pure-Python dependencies; UTF-8 output reconfiguration for Windows \\
\bottomrule
\end{tabular}
\caption{Non-functional requirements and how each is satisfied.}
\end{table}

\subsection*{3.10 Design Decisions and Trade-offs}
Several design choices involved genuine trade-offs. Recording them here explains why the system looks the way it does.

\item \textbf{Model choice: } A   t r a n s f o r m e r   w o u l d   l i k e l y   r a i s e   a c c u r a c y   o n   a   l a r g e r   c o r p u s ,   b u t   a t   t h i s   c o r p u s   s i z e   i t   w o u l d   o v e r f i t ,   w o u l d   a d d   h e a v y   d e p e n d e n c i e s ,   a n d   w o u l d   m a k e   d e c i s i o n s   h a r d e r   t o   j u s t i f y .   T h e   l i n e a r   m o d e l   w a s   c h o s e n   a s   t h e   d e f a u l t   w i t h   t h e   t r a n s f o r m e r   p a t h   r e t a i n e d   b e h i n d   a   f l a g .

\item \textbf{Persistence timing: } G u a r a n t e e s   c o m p l e t e   r e c o r d s   a n d   i m m e d i a t e   s u r f a c i n g   o f   f a i l u r e s ,   a t   t h e   c o s t   o f   a d d i n g   a n a l y s i s   t i m e   t o   s u b m i s s i o n   l a t e n c y .

\item \textbf{Fraud independence: } P r e v e n t s   a   l e g i t i m a t e   r e - f i l e   f r o m   b e i n g   s c o r e d   a s   b a d   f a i t h ,   a t   t h e   c o s t   o f   o n e   e x t r a   q u e r y   p e r   s u b m i s s i o n .

\item \textbf{Duplicate scoping: } T w o   i d e n t i c a l   c o m p l a i n t s   i n   d i f f e r e n t   w a r d s   a r e   a l m o s t   c e r t a i n l y   t w o   r e a l   i s s u e s ;   s c o p i n g   t o   t h e   w a r d   a v o i d s   s u p p r e s s i n g   t h e m .   T h e   c o s t   i s   t h a t   c r o s s - w a r d   d u p l i c a t e s   a r e   n o t   c a u g h t .

\item \textbf{Summarisation: } T h e   d e t e r m i n i s t i c   t e m p l a t e   a l w a y s   w o r k s   a n d   i s   r e p r o d u c i b l e ;   t h e   L L M   p a t h   i s   o p t - i n .   T h e   c o s t   i s   l e s s   f l u e n t   p r o s e .

\item \textbf{Storage: } Z e r o - c o n f i g u r a t i o n   a n d   p o r t a b l e   f o r   a   d e m o n s t r a t i o n ;   e x p l i c i t l y   n o t   a   p r o d u c t i o n   c h o i c e .

\item \textbf{Front end: } E l i m i n a t e s   J a v a S c r i p t   d e p e n d e n c i e s   a n d   g u a r a n t e e s   o f f l i n e   o p e r a t i o n ,   a t   t h e   c o s t   o f   i n t e r a c t i v i t y .

\subsection*{3.11 End-to-End Walkthrough of a Submission}
The following traces a single complaint through every stage, using the concrete example “Sewage overflowing near the market area” filed in the ward Wakad from the Mobile App channel.

\begin{table}[htbp]
\centering
\small
\begin{tabular}{@{}lll@{}}
\toprule
Step & What happens & Result for this complaint \\
\midrule
1. Intake & The form or API validates the payload and calls run\_pipeline() & description and ward captured \\
2. Preprocessing & Text is cleaned, tokenised, lemmatised and stop-words removed & normalised text produced \\
3. Duplicate check & Similarity is compared against prior Wakad complaints & no match above threshold \\
4. Classification & TF-IDF features feed the LinearSVC model & Public Safety (confidence about 0.50) \\
5. Fraud scoring & Similarity, frequency and pattern signals are combined & Low risk, score 0.0 \\
6. Clustering & The complaint is compared to open incident clusters & assigned to a cluster identifier \\
7. Priority & Six weighted signals are summed & High (score above 3.2) \\
8. XAI & The breakdown is converted to sentences & hazard wording and category risk cited \\
9. Persistence & The enriched record is written to the database & one complete row \\
10. Delivery & Dashboard KPIs, charts and tables recompute & queues and cards update \\
\bottomrule
\end{tabular}
\caption{End-to-end walkthrough of one submission through the eight-stage pipeline.}
\end{table}

This trace is also the mental model the smoke suite encodes: each numbered step corresponds to at least one automated assertion, so a regression in any stage fails the suite rather than silently degrading a dashboard number.

\subsection*{3.12 Security, Privacy and Ethical Design}
Although the project is a demonstration, several design choices were made with real-world deployment in mind:

\item \textbf{Data minimisation: } T h e   r e p o s i t o r y   c o n t a i n s   o n l y   s y n t h e t i c   c o m p l a i n t s   g e n e r a t e d   b y   t h e   s e e d   s c r i p t .   T h e r e   i s   n o   p e r s o n a l   i n f o r m a t i o n   t o   l e a k   a n d   n o   c o n s e n t   o b l i g a t i o n   t o   d i s c h a r g e .

\item \textbf{Injection safety: } C o m p l a i n t   t e x t   i s   t r e a t e d   a s   d a t a ,   n e v e r   e x e c u t e d ;   d a t a b a s e   a c c e s s   u s e s   p a r a m e t e r i s e d   q u e r i e s   t h r o u g h o u t ,   w h i c h   r e m o v e s   t h e   S Q L - i n j e c t i o n   c l a s s   o f   v u l n e r a b i l i t y .

\item \textbf{No automated punishment: } T h e   f r a u d   e n g i n e   p r o d u c e s   a   r i s k   s c o r e   a n d   a n   e x p l a n a t i o n   f o r   h u m a n   r e v i e w .   I t   n e v e r   b l o c k s   a   c i t i z e n ,   n e v e r   d e l e t e s   a   r e c o r d   a n d   n e v e r   m a r k s   a   p e r s o n   a s   f r a u d u l e n t   a u t o m a t i c a l l y .

\item \textbf{Duplicates are not fraud: } A   d u p l i c a t e   i s   r e c o r d e d   a s   a   d u p l i c a t e ,   n o t   a s   a n   o f f e n c e ;   t h e   s a m e   t e x t   c a n   b e   a   r e - f i l e   o r   a   f l o o d ,   a n d   o n l y   t h e   p a t t e r n   a c r o s s   m a n y   s u b m i s s i o n s   r a i s e s   t h e   s c o r e .

\item \textbf{Contestability: } E v e r y   s c o r e   c a r r i e s   t h e   r e a s o n s   t h a t   p r o d u c e d   i t ,   s o   a   c i t i z e n   o r   a u d i t o r   c a n   c h a l l e n g e   a   d e c i s i o n   o n   i t s   m e r i t s .

\item \textbf{Representational bias: } T h e   c l a s s i f i e r   l e a r n s   f r o m   a   c u r a t e d   c o r p u s ;   i f   t h a t   c o r p u s   u n d e r - r e p r e s e n t s   a   l o c a l i t y   o r   a   c o m p l a i n t   s t y l e ,   a c c u r a c y   w i l l   b e   u n e v e n .   T h i s   i s   t r e a t e d   a s   a   k n o w n   r i s k   i n   S e c t i o n   8   r a t h e r   t h a n   b e i n g   h i d d e n .

Authentication, role-based authorisation and audit logging are deliberately absent from this version and are listed as future work. They are prerequisites for any real deployment, not optional refinements.

\end{document}